\section{Introduction}\label{sec:introduction}

Local spin updates can preserve information about a realized equilibrium
configuration for a long time, even though the unconditioned chain is
stationary. We quantify this persistence throughout the low-temperature
phase of the zero-field Sherrington--Kirkpatrick model. The central issue
is the mass of slowly mixing initial states: a worst-state lower bound
alone gives no information about whether a Gibbs sample encounters the
obstruction.

For an integer $n\ge2$, let $\Sigma_n=\Sig$, and let
$(J_{ij})_{1\le i<j\le n}$ be independent standard Gaussian variables.
At inverse temperature $\beta>0$, the zero-field
Sherrington--Kirkpatrick model has Hamiltonian and Gibbs measure
\begin{equation}\label{eq:model}
 H^J(x)=H_n^J(x)=\frac1{\sqrt n}\sum_{i<j}J_{ij}x_ix_j,
 \qquad
 \pi^J(x)=\pi_{n,\beta}^J(x)
       =\frac{e^{\beta H^J(x)}}{Z_{n,\beta}^J},
 \qquad x\in\Sigma_n,
\end{equation}
where $Z^J=Z_{n,\beta}^J=\sum_{x\in\Sigma_n}e^{\beta H^J(x)}$.
For $x,y\in\Sigma_n$, their overlap is
\[
 R(x,y)=\frac1n\sum_{i=1}^n x_i y_i.
\]
The measure $\pi^J$ is invariant under the global spin reversal
$x\mapsto-x$. We fix $\beta>1$ before sending $n$ to infinity.

In continuous-time heat-bath dynamics, every site has an independent
rate-one clock. At a ring of site $i$, its spin is resampled according to
\[
 \PP(x_i\gets s\mid x_{[n]\setminus\{i\}})
  =\frac{\exp\{\beta s h_i^J(x)\}}
         {2\cosh(\beta h_i^J(x))},
 \qquad
 h_i^J(x)=\frac1{\sqrt n}\sum_{j\ne i}J_{ij}x_j,
 \quad s\in\{-1,1\},
\]
with $J_{ji}=J_{ij}$ and $J_{ii}=0$. Write $P_t^J$ for its transition
kernel. The total update-attempt rate is $n$, and an attempt may leave
the configuration unchanged. We also consider the discrete kernel $K^J$
that chooses one site uniformly and performs one such conditional
resampling. Both dynamics have stationary measure $\pi^J$.

For probability measures $\nu,\eta$ on $\Sigma_n$, set
$\norm{\nu-\eta}_{\TV}=\frac12\sum_x|\nu(x)-\eta(x)|$.
For a fixed realized state $x$, define
\begin{align}
 d^J(x,t)&=\norm{P_t^J(x,\cdot)-\pi^J}_{\TV},
 &T^J(x)&=\inf\{t\ge0:d^J(x,t)\le1/4\},
 \label{eq:point-mixing-time}\\
 d_{\mathrm{disc}}^J(x,k)&=
      \norm{(K^J)^k(x,\cdot)-\pi^J}_{\TV},
 &T_{\mathrm{disc}}^J(x)&=
      \min\{k\ge0:d_{\mathrm{disc}}^J(x,k)\le1/4\}.
 \label{eq:point-discrete-mixing-time}
\end{align}
Here $k$ is an integer, and neither mixing time takes a supremum over
the initial state.

\begin{theorem}\label{thm:main}
Set $\kap=1/10000$. For every fixed $\beta>1$,
\begin{align*}
 \PP_{J,\,x_0\sim\pi^J}
  \left\{d^J\bigl(x_0,\exp(n^\kap)\bigr)>\frac14\right\}
  &\longrightarrow1,\\
 \PP_{J,\,x_0\sim\pi^J}
  \left\{d_{\mathrm{disc}}^J
       \bigl(x_0,\lfloor\exp(n^\kap)\rfloor\bigr)>\frac14\right\}
  &\longrightarrow1.
\end{align*}
\end{theorem}

The order of conditioning in this statement is essential. We first draw
$x_0$ from $\pi^J$, then regard it as a fixed argument of the transition
kernel. Integrating that kernel over $x_0$ would instead recover
$\pi^J$ at every time. The joint formulation above is equivalent to
convergence of the corresponding Gibbs mass to one in disorder
probability. Indeed, if $F_n(J)\in[0,1]$ is the mass of the complementary
set, the joint assertion says $\EE F_n\to0$. Markov's inequality gives
convergence in probability, and boundedness proves the converse.

Total variation contracts under a Markov kernel. In particular, for
$s,t\ge0$,
\[
 d^J(x,t+s)
  =\norm{\bigl(P_t^J(x,\cdot)-\pi^J\bigr)P_s^J}_{\TV}
  \le d^J(x,t),
\]
and the discrete distances are nonincreasing in $k$ as well.
Since the state space is finite, $t\mapsto d^J(x,t)$ is continuous.
Consequently $d^J(x,t)>1/4$ implies $T^J(x)>t$; in discrete time,
$d_{\mathrm{disc}}^J(x,k)>1/4$ implies $T_{\mathrm{disc}}^J(x)>k$.
It follows that the theorem implies a lower bound at every fixed
polynomial time scale. Appendix~\ref{fixed-sec:crossings} will also give
an independent proof using a fixed number of crossing levels.

\begin{corollary}\label{cor:superpolynomial}
For every fixed $\beta>1$ and every fixed $a>0$,
\[
 \PP_{J,\,x_0\sim\pi^J}\{T^J(x_0)>n^a\}\longrightarrow1,
 \qquad
 \PP_{J,\,x_0\sim\pi^J}
       \{T_{\mathrm{disc}}^J(x_0)>n^a\}\longrightarrow1.
\]
Equivalently, each corresponding Gibbs mass tends to one in probability
over the disorder.
\end{corollary}

\begin{proof}
Fix $a>0$. Eventually $n^a$ is smaller than both
$\exp(n^\kap)$ and $\lfloor\exp(n^\kap)\rfloor$. Apply
Theorem~\ref{thm:main} at those larger times and use contraction of
total variation to pass to $n^a$.
\end{proof}

\subsection{Background and relation to earlier work}

Sherrington and Kirkpatrick introduced the model as a mean-field
description of a spin glass~\cite{SK75}. Parisi's variational
description of its equilibrium free energy~\cite{Parisi80} was
established rigorously through Guerra's interpolation bound and
Talagrand's matching lower bound for the limiting
pressure~\cite{Guerra03,Talagrand06}.
These equilibrium results provide the scalar variational structure
behind our argument. Passing from free energy to relaxation requires
additional information about which overlap configurations a local
trajectory can visit.

Ben Arous and Jagannath developed criteria that turn overlap
free-energy barriers into exponentially small spectral gaps for
mean-field spin glasses~\cite{BAJ18}. Their use of constrained replica
pressures is an important antecedent of the overlap comparisons used
here. Their barrier criteria require additional landscape hypotheses;
the explicit Ising sufficient conditions in that work exclude the
pure quadratic covariance of SK. They therefore do not directly give
the present result throughout $\beta>1$.
For SK itself, Sellke proves exponentially slow worst-case Glauber
mixing at sufficiently large inverse
temperature~\cite[Theorem~1.1]{Sellke25}.
That conclusion concerns the slowest initial states. It does not
determine the Gibbs mass of the states that remain far from equilibrium
after their initial configuration has been fixed.

The role of the external field is also substantial. Bandeira,
El Alaoui, and R\"odder prove polynomial worst-case mixing for lazy
discrete heat-bath dynamics at every fixed inverse temperature when
a sufficiently large constant uniform field is present, with high
probability over the disorder~\cite[Corollary~1.2]{BER26}.
Our result concerns zero field. Its distinctive combination is the
entire range of fixed $\beta>1$, Gibbs mass tending to one for
realized initial states, and a stretched-exponential time scale.
The exponent $1/10000$ is not asserted to be optimal, and the constants
in the proof may depend on $\beta$; the argument does not give
uniformity as $\beta$ decreases to one.

El Alaoui, Montanari and Sellke prove a complementary obstruction to
sampling at $\beta>1$ for algorithms stable under disorder perturbations,
in normalized $2$-Wasserstein distance~\cite[Theorem~2.6]{AMS22}. Their stability
hypothesis and approximation metric differ from the conditional Glauber
transition law considered here.

Numerical studies have examined much sharper predictions for
low-temperature relaxation. Billoire studies equilibrium
autocorrelation times for binary couplings and sequential Metropolis
updates~\cite{Billoire10}; Monthus and Garel compute inverse spectral
gaps for small SK systems with Metropolis
dynamics~\cite{MonthusGarel09}. Both report evidence associated with
an $n^{1/3}$ scale for logarithmic relaxation times. Their observables
and dynamics differ from the state-conditioned total variation
distance studied here, and these numerical predictions do not fix
the exponent in Theorem~\ref{thm:main}.

The static comparisons build on Guerra's Gaussian interpolation
framework~\cite{Guerra03} and Talagrand's extensions to systems with
prescribed replica overlaps~\cite[Propositions~3.4 and~4.4]{Talagrand07}.
Panchenko developed a general vector-spin formulation and proved the
asymptotic sharpness of these constrained-overlap
bounds~\cite{Panchenko18}. Talagrand's two- and three-copy bounds already
retain an explicit cost when the prescribed overlaps differ from the
covariance endpoint. In Section~\ref{static-sec:gaussian}, we derive the
finite-hierarchy matrix form needed here for the quadratic SK covariance,
allowing arbitrary finite replica dimension and nonunit endpoint
diagonals. The same section obtains a
quantitative scalar free-energy comparison by choosing an interpolation
path through a differential equation driven by overlap expectations.
This choice has a methodological parallel in the adaptive
interpolation method of Barbier and Macris~\cite{BarbierMacris19};
their theorems are not inputs to our finite-size SK estimate.
Gaussian concentration and the Brascamp--Lieb variance
inequality~\cite{BL76} supply the probabilistic estimates used in
these comparisons.

\subsection{Proof strategy}

Set $L=n^\kap$ and $\rho=L^{-1}$. Our task is to exclude an overlap
sign passage during the first $e^L$ units of time. To connect that task
to mixing, choose a reference $v_1$ with $R(v_1,x_0)\ge q_*$ for a
fixed $q_*>0$. The half-space $\{z:R(v_1,z)\le0\}$ has Gibbs mass at
least $1/2$ by spin reversal. A transition law within $1/4$ of
equilibrium must therefore enter it with probability at least $1/4$.
The reference will be chosen from independent equilibrium samples,
using the initial state but no future updates.

The first part of the paper establishes the static estimates that
make this path argument possible. Section~\ref{static-sec:gaussian}
compares the expected normalized log partition function
$n^{-1}\EE\log Z^J$ with the minimum of the scalar Parisi functional,
with error $O_\beta(n^{-1/24})$, and proves the matrix comparison
for constrained replica overlaps.
Section~\ref{static-sec:scalar} develops the scalar evolution and
the variational identities of a minimizing measure.
Section~\ref{static-sec:locking} uses these identities to compare
three replicas: when one pair has a substantial positive overlap, their
small signed overlaps with the third replica are forced to remain close,
apart from a set of very small averaged Gibbs probability.
Section~\ref{sec:energy} proves an equilibrium energy bound and a
quantitative estimate for small Gaussian rotations of the disorder.
All of these specialized inputs are proved within the paper.

There are two further static steps. First, the energy and rotation
bounds imply a coverage statement in Section~\ref{sec:coverage}:
simultaneously for every subset $S$, outside a set of targets of
absolute Gibbs mass less than $e^{-D_1L}$, each $x\in S$ has
Gibbs mass at least $e^{-D_3L}$ of $z\in S$ with $|R(x,z)|\ge q_*$.
Here $D_3$ is fixed after $D_1$. Sampling $\exp(O(L))$ independent
Gibbs configurations can therefore supply a reference inside $S$.
Second, Section~\ref{sec:narrow} treats triples whose two smaller
overlaps lie near a positive level $r$, while the larger overlap exceeds
$r$ by only a shrinking amount. It supplies locking at the scales
required by the backwards searches. The estimate uses a scalar
stationarity identity on a fixed interval; if that identity fails, a
pair-overlap pressure gap supplies a forbidden interval instead. This
alternative avoids assuming a complete description of the minimizing
measure's support.

Section~\ref{sec:dynamics} applies these estimates to a stationary
discrete path $(x_k)$, and obtains the continuous-time path by an
independent rate-$n$ Poisson clock. At a crossing of a small positive
overlap level, coverage supplies a new reference with large overlap
with the current configuration. Searching backwards for the new
reference's crossing, the locking estimates keep the overlaps with
previously retained references almost unchanged. Iterating creates
many simultaneous overlap constraints at a single path index.

The reference samples are arranged in blocks of a fixed length $B$.
Within a block, each new reference is chosen in the set of
configurations having small absolute overlap with all the earlier
references in that block. At the start of the next block this
restriction is reset. A conflict with older retained references can
then delete only a suffix of length at most a fixed $H_0<B$.
Since $H_0<B$, deletion cannot reach the first reference of the
preceding block or any older retained reference. In particular,
$v_1$ survives; Figure~\ref{fig:protected-block} shows this invariant.
This construction permits $p$ levels of order $L$
with only $\exp(O(L))$ total reference samples, while retaining a
fixed positive fraction of the levels.

The final probability estimate uses the order in which randomness is
revealed. At level $i$, the reference is already determined when the
independent threshold $r_i$ is drawn uniformly from an interval of
length $\rho$. Its overlap at any fixed path index is unlikely to fall
in the much narrower permitted window. Removing constraints by
conditional expectation in decreasing order of level bounds their
joint probability, even though later reference choices depend on
earlier thresholds. We sum over deterministic path indices and
deterministic subsets of retained levels. This avoids conditioning on
the adaptively chosen retained list and completes both clock versions.

Two auxiliary conclusions have a separate role. Appendix~\ref{poly-sec:coverage}
proves that every disorder-selected set of Gibbs mass at least
$n^{-D_1}$ has polynomially large conditional mass of pairs with
absolute overlap at least a fixed $q_*(\beta)$; this threshold is
chosen before $D_1$, and the resulting exponent $D_2$ is fixed before
$n$ grows. Its Gaussian-gain lemma and cell corollary are proved in
full. Appendix~\ref{fixed-sec:crossings} combines that result with the
fixed-scale locking estimate to obtain a second proof of the
polynomial-time obstruction. There the number of levels is fixed and
the reference banks have polynomial sizes. These appendices reuse the
energy, rotation and static locking proofs already established in the
body, while retaining their own complete coverage and crossing arguments.

\paragraph{Notation.}
We suppress $J$, $n$, and the fixed $\beta$ from $H_n^J$, $\pi_{n,\beta}^J$,
and $Z_{n,\beta}^J$ when harmless. All constants may depend on $\beta$
and on other parameters fixed before $n$ grows. Gibbs replicas and
reference samples are independent conditional on the disorder unless
stated otherwise. For matrices, $U:V=\sum_{a,b}U_{ab}V_{ab}$ and
$|U|^2=U:U$ denote the Frobenius inner product and squared norm.
